\section{Intermediate terminal orders}\label{app:middle-tail}

\begin{proposition}\label{prop:middle-tail}
Assume the conditions of Lemma~\ref{lem:support-reduction}.
For \(2200\le N<e^{15}\), a common factor requires
\(\alpha<.03\) and \(b>N\). For \(3150\le N<e^{15}\)
there is no common factor.
\end{proposition}

Write \(t=\log(1/\alpha)\), \(V=\log(M+1)\) and
\(\sigma=2U+2\), where \(U=\log N\).
Throughout this appendix,
\(\log m=U+\log(1-1/N)>U-1/(N-1)>U-.001\).
From Lemma~\ref{lem:gap-caps} we may use \(d\le233\) below
\(20{,}000\) and \(d\le565\) throughout the appendix.
If \(X\) bounds the prime divisors of the lifted norm \(K\), the
right side of \eqref{eq:pade-norm} is at most
\begin{equation}\label{eq:intermediate-prime-head}
 S(X)+\pi(X)V/N,\qquad S(X)=\sum_{p\le X}\frac{\log p}{p-1}.
\end{equation}
In the applications below \(M\le X^j\), \(M\ge7N\) and
\(X\ge9\) give
\[
 \frac{V}{\log X}
 \le j+\frac{\log(1+1/M)}{\log X}
 <j+\frac{1}{7N\log9}<j+.001 .
\]
The baseline from Lemma~\ref{lem:resultant} is
\(\log\Lambda> U-1.011-g\). The weaker constant \(1.016\)
also suffices for the estimates of proper cosets.

\begin{lemma}\label{lem:small-prime-sums}
For each of the pairs
\[
 (x_0,c)=(21,.22),(34,.30),(43,.33),(100,.39),
          (117,.425),(137,.449),(150,.45),(343,.49)
\]
one has \(S(x)<\log x-c\) for \(x\ge x_0\).
For the prime count one has
\[
 \frac{\pi(x)\log x}{x}<
 \begin{cases}
  1.256 & (x\ge40),\\
  1.224 & (x\ge117),
 \end{cases}
\]
with the smaller bound on the overlap.
\end{lemma}

\begin{proof}
For \(S(x)-\log x\) below \(3000\), check the listed left endpoints
and the primes, since the function decreases between primes.
For the prime count, check \(40\) and \(117\) and the primes
up to \(300\), since between primes the multiplier \(\log x/x\)
decreases for \(x>e\). For \(x\ge300\),
Lemma~\ref{lem:prime-estimates} gives
\(1+1.2762/\log x\le1+1.2762/\log300<1.224\).
The supplement encloses the logarithms at all these checkpoints.
\end{proof}

\subsection{Proper rotation cosets}

Let \(k_0\) be the size of each rotation coset in
Lemma~\ref{lem:kummer}. Suppose \(q=d/k_0\ge2\), and put
\(P=M^{1/(q+1)}\). Then all norm primes are at most \(P\)
and \(4d\le2H<P\). In particular,
\begin{equation}\label{eq:proper-coset-two-caps}
 P^q<\min\{N/(4\alpha),N^2/4\}.
\end{equation}
The first inequality uses \(M=dN/\alpha<PN/(4\alpha)\),
and the second uses \(M<dN^2\).

If \(k_0=1\) and \(d\ge3\), then \(P^{d+1}=M<dN^2\).
Together with \(4d<P\) this gives \(4(4d)^d<N^2\).
In particular \(d\log(4d)<2U\), so \(d<U\) because
\(\log(4d)\ge\log12>2\).
Use \(X=(UN^2)^{1/4}>P\) and \(s<3U/4\) in the
baseline with \(g\le7s/32\). The resulting surplus over
\eqref{eq:intermediate-prime-head} is at least
\[
 U-1.016-\frac{21}{128}U
 -\log X+.22-4.001\cdot1.256\,X/N.
\]
Here \(\log X=U/2+(\log U)/4\).
The derivative of the terms linear or logarithmic in \(U\) is
\(1-21/128-1/2-1/(4U)>0\), and \(X/N\) decreases because
\(-1/2+1/(4U)<0\). The surplus therefore exceeds the certified
value \(1.10\) at \(N=2200\).
If \(k_0=1,d=2\), take \(X=(2N^2)^{1/3}\) and \(s<1.5\).
Now \(\log X=(2U+\log2)/3\), so the surplus again increases.
With the prime-sum correction \(.33\) in place of \(.22\),
it exceeds \(.95\) at \(N=2200\).

Suppose \(k_0\ge2\) and \(\alpha\ge.03\). For
\(\alpha\in[\ell,u]\) put \(X=\sqrt{N/(4\ell)}\).
By \eqref{eq:proper-coset-two-caps}, \(X>P\) and \(M<X^3\).
At least a third of each coset lies in the left half-plane.
The one-sided gain from Lemma~\ref{lem:resultant} gives the
lower bound
\[
 U-.006-\tfrac23(1.01+\kappa\sigma)
       +\tfrac13\log((1-u)/u)
\]
on \(\log\Lambda\), using \(g\le\kappa s\).
In every row \(X>40\), so \eqref{eq:intermediate-prime-head}
is bounded above by \(\log X+c+3.001\cdot1.256\,X/N\).
The lower bound on the gain has slope \(1-4\kappa/3>1/2\)
in \(U\). The logarithmic part \(\log X\) of the upper bound
has slope \(1/2\), and \(X/N\) decreases. It remains to check
the following margins at \(N=2200\).
\[
\begin{array}{cccc}
\toprule
{[\ell,u]}&\kappa&c&\text{surplus exceeds}\\
\midrule
{[.03,.1]}&.134&-.30&1.35\\
{[.1,1/3]}&.160&.20&.75\\
{[1/3,.75]}&.218&.20&.14\\
\bottomrule
\end{array}
\]

For \(\alpha<.03\) and \(q\ge3\), take
\(X=N^{2/3}/4^{1/3}>P\). From \(4d<P\) we obtain
\(M<dN^2<XN^2/4=X^4\).
If \(D\) bounds \(d\), the baseline exceeds
\(U-1.016-.00384D\).
Subtracting the upper bound
\(\log X-.22+4.001(X/N)(1+1.2762/\log X)\)
gives a surplus that increases for fixed \(D\), since
\(\log X\) has slope \(2/3\) in \(U\) and
\((X/N)(1+1.2762/\log X)\) decreases.
At \(N=2200,D=233\) and \(N=20{,}000,D=565\) it
exceeds \(1.08\) and \(.68\), respectively.

For \(\alpha<.03\) and \(q=2\), use \(X=N/2\), \(M<X^3\)
and \eqref{eq:intermediate-prime-head} to bound the norm
charge by \(U+.862\).
When \(d=4\), half of the size-two coset has gain at least
\(m(1-\alpha)/\alpha\), and the mean log gain exceeds
\(U+1.21\). For \(6\le d\le18\), the small-coset form of
Lemma~\ref{lem:angles}, with \(k_0=d/2\), gives a gain greater
than \(U+1.89\). For \(d\ge20\), use \(g_0\le3.18\),
\(g_0/k_0\le.10868\), \(r/k_0<.111\) and \(J<.153\)
in \eqref{eq:pade-angles}, together with
\(\log C'>t-.7237\) and \(1-1/k_0\ge.9\).
Dropping the positive \(2\log(2k_0/\pi)/k_0\) leaves a gain
exceeding \(U+1.80\).
This proves that the rotation coset is full.

\subsection{Full cosets at moderate and large ratios}

First take \(2200\le N<20{,}000\).
Here \(a\le(97/3)d\), \(d\le233\) and \(H\le467\).
Thus \(aN+H<189{,}695{,}659=X_{222}\).
Lemma~\ref{lem:local-gap} gives \(H<222\) and \(d\le110\).
For \(N\le3500\) these caps give \(aN+H<X_{154}=17{,}051{,}707\),
so \(d\le76\).

Let \(\ell\le\alpha\le u\) and order the \(d\) rotations by
increasing absolute argument, starting at index zero. An arc
centered at angle zero and containing \(j+1\) arguments has
length at least \(2\pi j/d\), by the rotation spacing.
Thus \(c_j=\cos(\pi j/d)\) bounds the cosine in position \(j\).

For a nonnegative upper bound \(c\) on the cosine, set
\(W=(1-u-c)/u\). In the second integrand in
\eqref{eq:pade-integrals}, write \(v=1-e^{-x/d}\).
Use \(v\le x/d\) and \(\Re z/m\le bc\) to bound the logarithm
of the original integrand, divided by \(m\), by
\[
 -\frac{1-\alpha-c}{\alpha}x+\log(1-e^{-x})
 \le -Wx+\log(1-e^{-x}).
\]
For \(W>0\) the derivative on the right vanishes only at
\(e^{-x}=W/(1+W)\), and the second derivative is negative.
The maximum is \(-J(W)\), where
\(J(W)=(1+W)\log(1+W)-W\log W\).
The original integral over \(0\le v\le1\) is at most
\(e^{-mJ(W)}\). Since \(J(W)<1+\log(1+W)\le
1+\log(100/3)<5\), Stirling's inequality gives
\((m!e^{mJ(W)})^{1/(m+1)}>me^{-1.01+J(W)}\).
Indeed \(.01(m+1)+\tfrac12\log(2\pi/m)-5\) is positive at
\(m=2199\) and has derivative \(.01-1/(2m)>0\).
The linear gain in Lemma~\ref{lem:resultant} is at least \(mW\).
Together the two gains give
\(\Phi_u(c)=\max\{-1.01+J(W),\log W\}\) after subtracting
\(\log m\).
For \(c<0\) use \(W=(1-u)/u\) in the integrand and
replace \(\log W\) in this maximum by \(\log(W(1-c))\).
For \(c\ge0,W\le0\) use the baseline
\(-1.01-\kappa(u)du\), where \(\kappa(u)=g/s\) at \(\alpha=u\)
in \eqref{eq:pade-loss}.
In the smaller terminal band we also use the lower bound
\[
 -1.01-\frac d u
       \left(c-(1-u)-(1-u)\log\frac{c}{1-u}\right)
\]
when this improves the baseline. It follows by discarding
the beta tail factor and maximizing the remaining exponent.
The loss at \(\alpha\) is zero for \(c\le1-\alpha\) and
otherwise the same formula with \(\alpha\) in place of \(u\).
For \(c>1-\alpha\) the derivative of the loss is
\[
 \frac d{\alpha^2}
 \left(\log\frac{c}{1-\alpha}-c+1-\alpha\right)\ge0.
\]
For the inequality, \(\log(c/(1-\alpha))\) is the integral
of \(1/x\ge1\) from \(1-\alpha\) to \(c\le1\).
This permits evaluating the loss at the upper endpoint \(u\).

By Lemma~\ref{lem:degree-cap}, a prime bound is
\(\max(a+2d-1,\sqrt M)\). It lies between \(X_1\) and \(X_0\),
where
\[
 X_0=\max\{d(1+1/\ell),\sqrt{dN/\ell}\},\qquad
 X_1=\sqrt{\max(7,d/u)N}.
\]
The upper bound for the right side of the norm inequality is therefore
\begin{equation}\label{eq:moderate-norm}
 \log X_0-c_0+2.001(1+1.2762/\log X_1)X_0/N.
\end{equation}
Use \(c_0=.39\) in the first band, where \(X_1>100\), and
\(c_0=.45\) in the second, where \(X_1>150\).
In \(U\), \(\log X_0\) is the maximum of a constant and an
affine function of slope \(1/2\). The quotient \(X_0/N\) is
the maximum of two decreasing functions, and \(1/\log X_1\)
also decreases. Thus the norm bound grows at most at rate
\(1/2\). The angular lower bound
\(\log(N-1)+d^{-1}\sum_j\Phi_u(c_j)\) grows at rate \(N/(N-1)>1\).
Since \(d\ell\le s<2U+2\), the first eligible order is
bounded below by \(\max\{N_{\min},\exp((d\ell-2)/2)\}\).
For each \(d\) and each interval of the following subdivision,
the interval checker evaluates the surplus at that lower bound.
Intervals with \(d(1/\ell-1)<5\) contain no eligible support.
\[
\begin{array}{rrlc}
\toprule
N_{\min}&d_{\max}&\text{interval widths}&\text{surplus exceeds}\\
\midrule
2200 &76 & .005\text{ to } .05,\ \text{then } .01 &.070\\
3500 &110& .02 &.089\\
\bottomrule
\end{array}
\]

Now take \(20{,}000\le N<e^{15}\). The gap table gives the
successive caps \(D=257,300,565\) at starting orders
\(20{,}000,32{,}768,65{,}536\). In each range \(D/N<.015\).
Since \((1+\alpha)^2/\alpha\) decreases on \([.03,.75]\),
\[
 \frac{(1+\alpha)^2b}{N}
 \le\frac D N\frac{(1+\alpha)^2}{\alpha}
 <.015\frac{1.03^2}{.03}<1.
\]
Hence \(\sqrt M\) bounds the norm primes. Using
\(S(\sqrt M)<\log\sqrt M-.49\), the right side of the norm inequality is below
\begin{equation}\label{eq:moderate-high-norm}
 \tfrac12(U+\log b)-.49+2.46\sqrt{b/N}.
\end{equation}
At angles with \(\cos\theta\le1-\alpha\) the beta integral
has gain at least \(me^{-1.01}\). The exceptional fraction
is at most \(\arccos(1-\alpha)/\pi+1/d\). For
\(\ell\le\alpha\le u\), use \(\eta\ge\arccos(1-u)/\pi\).
At \(d\ge10\), at least a fraction \(.45\) lies in the left
half-plane, where one may use the greater gain \(m(1-\alpha)/\alpha\).
Thus, with \(S_0=\min(2U+2,Du)\) and
\(b_0=\min((2U+2)/\ell^2,D/\ell)\), the surplus exceeds
\[
 U-1.011-\kappa(\eta S_0+u)
  +.45\max\{0,1.01+\log((1-u)/u)\}
  -\bigl(\tfrac12(U+\log b_0)-.49+2.46\sqrt{b_0/N}\bigr).
\]
Where \(S_0\) is differentiable in \(U\), its derivative is
zero or two. Thus the lower bound on the gain grows at rate
at least \(1-2\kappa\eta\). The logarithmic term in the
upper bound grows at most at rate \(1/2+1/(2U+2)\), since
\((\log b_0)'\le2/(2U+2)<1\). The latter inequality also
makes \(\sqrt{b_0/N}\) decrease. The checker verifies
\(1-2\kappa\eta>1/2+1/(2U+2)\) at each starting order,
where the comparison is tightest. It also checks
\(\cos(\pi\eta)<1-u\) and the upper bound \(\kappa\)
on \(g/s\) at \(u\). The table records the least certified
surplus across the three starting orders for each interval.
For \(d\le9\) the same argument before adding the left half-plane
gain uses \(\kappa=.218,\eta=.42,\ s\le6.75,\ b\le300\),
and gives a positive surplus at \(N=20{,}000\) increasing thereafter.
\[
\begin{array}{cccc}
\toprule
{[\ell,u]}&\kappa&\eta&\text{surplus exceeds}\\
\midrule
{[.03,.036]}&.128&.086&.12\\
{[.036,.045]}&.129&.096&.20\\
{[.045,.07]}&.131&.120&.11\\
{[.07,.11]}&.134&.151&.21\\
{[.11,.17]}&.140&.189&.51\\
{[.17,.28]}&.150&.245&.67\\
{[.28,.45]}&.170&.315&.62\\
{[.45,.60]}&.190&.370&.53\\
{[.60,.75]}&.218&.420&.08\\
\bottomrule
\end{array}
\]


\subsection{Small ratios with the upper exponent at most the order}

Take \(X=1.03\sqrt M\), which is a prime bound by
Lemma~\ref{lem:degree-cap}. Indeed
\(a+2d-1<(1+\alpha)b\le1.03\sqrt M\).
For \(d\ge20\), the angular estimate
\eqref{eq:pade-angles} gives
\[
 \log\Lambda-U>
 .95\bigl(\log(100/3)-.7237\bigr)-.002-.05434
 - .111\bigl(1.01+.00384D+\log2\bigr)-.153
\]
with \(d\le D\). In \eqref{eq:intermediate-prime-head},
use \(X/N\le1.03\) and \(X>1.03\sqrt{20N/.03}\).
At \(N=2200,D=233\) the surplus exceeds \(.17\), and at
\(N=20{,}000,D=565\) it exceeds \(.08\).
For fixed \(D\), the gain bound minus \(U\) is constant.
After using \(X/N\le1.03\) in the norm bound, its remaining
\(N\)-dependence is through the decreasing reciprocal of
\(\log(1.03\sqrt{20N/.03})\).
For \(10\le d\le19\), inserting each \(d\) in
\eqref{eq:pade-angles} gives \(\log\Lambda>U+2.41\),
against a right side of the norm inequality below \(U+2.052\) by \(X>343\).
For \(4\le d\le9\) use the small-coset angular bound. The gain
exceeds \(U+2.24\), against a right side of the norm inequality below \(U+2.177\)
by \(X>137\).

For an antipodal coset with \(d=2\), write \(c=|\cos\theta|\),
\(A_0=(1-\alpha)/\alpha\) and \(B_0=e^{-g_0}\).
The two gains have product at least
\(m^2A_0(1+c)\max\{B_0,A_0-c/\alpha\}\).
On \(0\le c\le1-\alpha(1+B_0)\) this is a concave quadratic
in \(c\) and is linear increasing on the rest of its domain.
The two endpoint values in the quadratic range, divided by
\(m^2A_0\), are \(A_0\) and \(B_0(2-\alpha(1+B_0))\).
The second is smaller because \(A_0>32\) and \(B_0<1\).
This gives the lower bound \(m^2A_0B_0(2-\alpha(1+B_0))\).
In particular \(\log\Lambda>U+t/2-.2\).
For \(d=3\) the small-coset bound gives
\(U+2t/3-.435>U+t/2-.2\).
Put \(\mu=\tfrac12\log(b/N)\).
Since \(X>117\), the right side of the norm inequality is at most
\[
 U+\mu+\log(1.03)-.425+
       2.001\cdot1.03\cdot1.224\,e^\mu.
\]
Substitute \(t=U-\log d+2\mu\) in the gain \(U+t/2-.2\).
The linear terms in \(\mu\) cancel in the surplus, while the
remaining negative multiple of \(e^\mu\) is smallest at
\(\mu=0\). The resulting bound increases with \(U\)
and exceeds \(.96\) at \(N=2200\).
This proves the first conclusion of Proposition~\ref{prop:middle-tail}.

\subsection{Small ratios with the upper exponent exceeding the order}

It remains to exclude this case for \(N\ge3150\).
The bounds \(H<G,\ d\le D=G/2-1\) are shown in the
following table. In addition to the main gap table, apply
Lemma~\ref{lem:local-gap} successively using \(G=320,292,288,282\)
up to \(N=5000,4150,3467,3295\), respectively.
At every step the preceding caps give
\(aN+H<(5/6)D_{\mathrm{old}}N_{\max}^2+H_{\mathrm{old}}<X_G\).
\[
\begin{array}{rrrccc}
\toprule
N\text{ from}&N\text{ through}&G&A_v&A_c&\text{surplus exceeds}\\
\midrule
3150  &3295   &282 &2.504&2.885&.023\\
3296  &3467   &288 &2.506&2.886&.146\\
3468  &4150   &292 &2.514&2.889&.31\\
4151  &5000   &320 &2.522&2.892&.44\\
5001  &6000   &354 &2.531&2.896&.86\\
6001  &10000  &456 &2.58 &3.01 &.59\\
10001 &19999  &468 &2.58 &2.92 &2.37\\
20000 &32767  &516 &3&3&3.88\\
32768 &65535  &602 &3&3&4.66\\
65536 &3269017&1132&3&3&4.87\\
\bottomrule
\end{array}
\]

Let
\[
 \begin{aligned}
 w&=\log(b/N), & v&=\log M=2U+w, & P&=4M^{1/3},\\
 \zeta&=\log P, & t_D&=U+w-\log D, & \xi&=t_D+.337.
 \end{aligned}
\]
We have \(0<w<\log(5D/6)\) and \(P>\max(830,5G/2)\)
in every row. At \(p>P\) we have \(p>2H+d\) and \(p^3>M\).
Lemma~\ref{lem:two-heights} separates the prime divisors of the
norm there according to whether the two Newton edge heights agree.

The primes with equal heights divide
\(\binom{a+2d-1}{2d-1}\), since they divide a number in
\([a+1,a+2d-1]\) and exceed \(2d-1\).
Using \(\binom{x}{j}\le(ex/j)^j\), their total logarithmic weight
is at most
\[
 2d\bigl(1+\log(a/d+2)-\log2\bigr)<2d(t+.337)\le2D\xi.
\]
For the last step the function \(x(\log(b/x)+.337)\) increases
for \(2\le x\le D\), since \(b>N\) and \(D<N/4\).
Their right side of the norm inequality is at most
\(2D\xi(1/(P-1)+V/(N\zeta))\).

For unequal heights, one of the two windows
\([bm-2H,bm]\), \(m=N-1,N\), contains \(qp^2\)
with \(1\le q\le Q=M/P^2=P/64\).
At fixed \(q\) each window contains at most one such \(p>P\),
since consecutive integer squares multiplied by \(q\) are
spaced by more than \(2P>2H\) there.
Thus the digit charge is at most \(2QV/N\).
Put \(Z=b(N-1)-2H\). Since \(Z>.9995M\) and \(Q>12\),
the first sum in \eqref{eq:pade-norm} on these primes is at most
\[
 v\,\frac{\sum_{q\le Q}\sqrt q}{\sqrt Z-\sqrt Q}
 \le \frac{1.002v}{64}\left(\frac23+\frac{1}{2Q}\right)
 < .71v/64.
\]
Here \(p\ge\sqrt{Z/q}\) and \(p\le\sqrt M\).
For the middle bound use the trapezoid inequality for the concave
square root and
\(\sqrt Z-\sqrt Q\ge\sqrt M(.9995-1/P)>\sqrt M/1.002\).
The lower bound on \(Z/M\) follows from
\(1-1/N-2H/N^2>.9995\), with \(H\le467\) below
\(20{,}000\) and \(H\le1131\) above.

The primes \(p\le P\) contribute at most
\(S(P)+\pi(P)V/N\). We have \(V<v+.001\) and
\[
 \frac{V}{\zeta}<3-\frac{3\log4-.001}{\zeta}.
\]
Put \(k=3\log4-.001\) and \(c_1=1.2762\).
Both \(3-k/\zeta\) and
\[
 (3-k/\zeta)(1+c_1/\zeta)
 =3-\frac{k-3c_1}{\zeta}-\frac{kc_1}{\zeta^2}
\]
increase in \(\zeta\), since \(k>3c_1\).
This also bounds both by three. In each row the upper corner
has its largest \(N\) and \(b/N=5D/6\), so the checker
evaluates there to obtain \(A_v,A_c\).
Combining the three ranges of primes, using \(S(P)<\log P-.49\)
and \(P/(P-1)<1.002\), gives the upper bound
\begin{equation}\label{eq:hard-upper}
 R=\zeta-.49+\frac PN\left(A_c+\frac{2(v+.001)}{64}\right)
 +\frac{.71v}{64}+\frac{2.004D\xi}{P}+\frac{2A_vD\xi}{N}.
\end{equation}

For the lower bound let \(E=D/b\).
The errors in \eqref{eq:pade-angles} satisfy
\(J<.883\sqrt\alpha\) and \(r/d<.641\sqrt\alpha\).
Indeed the derivative of \(\arcsin x+x\sqrt{1-x^2}\)
is \(2\sqrt{1-x^2}\le2\). For the second estimate,
convexity bounds \((1-y)^{-1/2}\) on \(0\le y\le.03\)
by the chord from \(1\) to \(1/\sqrt{.97}<1.0156\).
Integrating the resulting upper bound \(1+.52x^2\) gives
\(\arcsin x\le x+.52x^3/3<1.006x\) for \(0<x^2\le.03\).
For \(d\ge20\), use
\(\log C'>t-.6932-1.031E\),
\(g_0\le1.01+.128DE\), and \(g_0/d\le1.01/d+.128E\).
At fixed \(b\), the derivatives of
\((1-1/x)(\log(b/x)-.6932)\) and
\([2\log(2x/\pi)-1.01]/x\), multiplied by \(x^2\), are
\[
 \log(b/x)-.6932-x+1,\qquad 3.01-2\log(2x/\pi),
\]
respectively. Both are negative for \(d\le x\le D\), since
\(x\ge20\) and \(\log(b/x)\le t<U<15\).
Substituting \(D\) for \(d\) in these two decreasing terms
and bounding the errors with \(E\) gives
\begin{equation}\label{eq:hard-lower}
 \begin{aligned}
 L={}&U-.001+(1-1/D)(t_D-.6932)
     +\frac{2\log(2D/\pi)-1.01}{D}\\
     &-1.159E-\bigl[.641(1.7032+.128DE)+.883\bigr]\sqrt E .
 \end{aligned}
\end{equation}
For \(d=2,3,\ 4\le d\le9,\ 10\le d\le19\), use
\(D=2,3,9,19\) respectively in \eqref{eq:hard-upper} and
\(t_D\), keeping the row's \(A_v,A_c\).
The same angular estimates used above give
\(L=U+ft_D-c\) with
\((f,c)=(1/2,.2),(2/3,.435),(.75,.68),(.9,.925)\),
respectively. For \(10\le d\le19\) retain the positive term
\(2\log(2d/\pi)/d\ge2\log(38/\pi)/19\) in
\eqref{eq:pade-angles}.

For fixed \(U\), the lower bound \(L\) is affine in \(w\)
minus a nonnegative linear combination of \(e^{-w}\),
\(e^{-w/2}\) and \(e^{-3w/2}\). It is therefore concave.
For \(R\) only the \(P/N\) and \(\xi/P\) terms have nonzero
second derivatives in \(w\). Their sum gives
\[
 R_{ww}=
 \frac{P}{9N}\left(A_c+\frac{2(v+.001+6)}{64}\right)
      +\frac{2.004D}{9P}(\xi-6).
\]
To prove \(R_{ww}\ge0\), multiply by \(9P\) and compare the
positive term with \(2.004D(6-\xi)\).
The positive term increases in \(U,w\) because
\(P^2/N=16e^{U/3+2w/3}\) and \(v=2U+w\).
The term \(D(6-\xi)\) decreases in \(U,w\).
Where it is positive, it increases with \(D\) because its
derivative in \(D\) is \(7-\xi>1\).
It suffices to check \(w=0\) at the smallest order in each row
with its largest \(D\).

For fixed \(w\), \(L\) grows with \(U\) at rate at least \(1.5\),
since its negative exponential terms increase towards zero.
The \(U\)-derivative of \((P/N)(A_c+(v+.001)/32)\) equals
\[
 \frac PN\left(\frac1{16}-\frac{A_c+(v+.001)/32}{3}\right)<0,
\]
since \(A_c>2\).
The two equal-height terms decrease because their derivatives
have the signs of \(1-2\xi/3\) and \(1-\xi\), and \(\xi>3\).
Thus \(R\) grows at rate at most \(2/3+2\cdot .71/64\).
The supplement checks the convexity inequality and the surplus
at the two endpoints of \(w\) at the smallest \(N\) in each row.
For \(d\ge20\) the certified margins are given in the table.
For all four smaller-gap cases the margin exceeds \(3.8\).
It follows that \(L>R\) throughout each row, which contradicts
\eqref{eq:pade-norm} and completes the proposition.

